% Closure of wp3 thm:PLDI  [v2.1 — self-contained: quadratic KL minorant,
% Hanson-Wright pointwise tails, and an explicit two-dimensional chaining
% lemma for the sub-exponential increments. Delivers an exponential bound,
% stronger than the stated polynomial one.]
\closuredraft
\begin{lemma}[Two-dimensional sub-exponential chaining]\label{lem:chain2d}
Let $\{X(u):u\in\mathcal B\}$ be a separable stochastic process indexed by a
ball $\mathcal B\subset\R^2$ of radius $R>0$, with $X(u_\ast)=0$ for some
$u_\ast\in\mathcal B$ and increments satisfying, for all
$u,u'\in\mathcal B$ and $t>0$,
\begin{equation}\label{eq:chain-inc}
 \Pp\big(|X(u)-X(u')|\ge t\big)
 \le2\exp\Big\{-c_0\min\Big(\frac{t^2}{L^2\norm{u-u'}^2},\
 \frac{t}{L\norm{u-u'}}\Big)\Big\}
\end{equation}
with constants $c_0,L>0$. Then there are $C_1,c_1>0$ depending only on
$c_0$ such that, for all $x\ge C_1$,
\[
 \Pp\Big(\sup_{u\in\mathcal B}|X(u)|\ge x\,L\,R\Big)
 \le C_1\exp\{-c_1x\}.
\]
\end{lemma}
\begin{proof}
For $j\ge0$ let $\mathcal N_j\subset\mathcal B$ be a $2^{-j}R$-net with
$\#\mathcal N_j\le C4^{\,j}$ (possible in $\R^2$), $\pi_j$ a
nearest-net-point map, $\pi_0\equiv u_\ast$. By separability it suffices to
bound the supremum over $\bigcup_j\mathcal N_j$, and along the chain
$X(u)=\sum_{j\ge1}\{X(\pi_ju)-X(\pi_{j-1}u)\}$ with
$\norm{\pi_ju-\pi_{j-1}u}\le3\cdot2^{-j}R$. Allocate
$t_j:=\eta\,xLR\,j2^{-j/2}$, $\eta:=(\sum_{j\ge1}j2^{-j/2})^{-1}$. At stage
$j$ there are at most $C\,16^{\,j}$ increment pairs, so by
\eqref{eq:chain-inc},
\[
 \Pp\Big(\sup_{\mathcal B}|X|\ge xLR\Big)
 \le\sum_{j\ge1}C\,16^{\,j}
 \exp\Big\{-c_0\min\Big(\tfrac{\eta^2}{9}x^2j^24^{\,j/2},\
 \tfrac{\eta}{3}x\,j\,2^{\,j/2}\Big)\Big\}
 \le\sum_{j\ge1}C\,16^{\,j}e^{-c\,x\,j\,2^{\,j/2}}
 \le C_1e^{-c_1x},
\]
where the middle inequality uses $x\ge1$ (so the second entry of the
minimum is active up to constants) and the last holds for $x\ge C_1$ with
$C_1$ large enough that $c\,x\,j\,2^{\,j/2}\ge2j\log16+c_1x(1+j)$ for all
$j\ge1$.
\end{proof}

\closuredraft
\begin{proof}[Proof of the theorem]
Throughout $u:=\An(\vartheta-\vzero)\in\R^2$, and $c,C$ depend only on
$\Theta,T,x_0$.

\emph{Step 1: global quadratic KL minorant.}
We claim that for $N,n$ large,
\begin{equation}\label{eq:pldi-minorant}
 K(\vartheta):=\E_{\vzero}\big[\elln(\vzero)-\elln(\vartheta)\big]
 \ \ge\ c\,\norm{u}^2
 \qquad\text{for all }\vartheta\in\Theta .
\end{equation}
Start from the exact decomposition \eqref{eq:KL-exact} in the proof of Lemma~\ref{lem:KL}. \emph{Drift part:}
by the mean value theorem bound in the proof of Lemma~\ref{lem:KL}, uniformly on $\Theta$,
$\text{drift}\ \ge c(\theta-\theta_0)^2\LamN=c\,u_\theta^2$. \emph{Scale
part:} write $v^0_k/v_k(\vartheta)=x\,(1+e_k)$ with $x:=\sigma_0^2/\sigma^2$
and $e_k:=h^0_k/h_{k,n}(\theta)-1$, $|e_k|\le C\lambda_k\Delta_n
|\theta-\theta_0|$ (Lipschitz in $\theta$ via \eqref{eq:g}). Since
$\psi'(y)=1-1/y$ satisfies $|\psi'(y)|\le C(|x-1|+|e_k|)$ for $y$ between
$x$ and $x(1+e_k)$ on the compact range,
\[
 \big|\psi\big(x(1+e_k)\big)-\psi(x)\big|
 \le C\big(|\sigma-\sigma_0|+\lambda_k\Delta_n|\theta-\theta_0|\big)
 \lambda_k\Delta_n|\theta-\theta_0| .
\]
Summing over $(k,j)$ and using $\psi(x)\ge c(\sigma-\sigma_0)^2$ on
compacts, $n\Delta_n\sum_k\lambda_k=T\LamN$ and
$n\Delta_n^2\sum_k\lambda_k^2\le T\lambda_N\Delta_n\LamN$:
\[
 \text{scale}\ \ge\ cNn(\sigma-\sigma_0)^2
 -CT\LamN|\sigma-\sigma_0||\theta-\theta_0|
 -CT(\lambda_N\Delta_n)\LamN(\theta-\theta_0)^2 .
\]
The last term is $o(1)\cdot u_\theta^2$. For the middle term, AM--GM gives
$\LamN|\sigma-\sigma_0||\theta-\theta_0|\le\tfrac{\epsilon}{2}\LamN
(\theta-\theta_0)^2+\tfrac{1}{2\epsilon}\LamN(\sigma-\sigma_0)^2$, and
$\LamN/(Nn)\le\lambda_N\Delta_n/T\to0$ makes the second piece
$o(1)\cdot Nn(\sigma-\sigma_0)^2$. Choosing $\epsilon$ small and combining
with the drift part yields \eqref{eq:pldi-minorant}.

\emph{Step 2: Gaussian-quadratic representation.}
Let $\xi=(\varepsilon_{k,j})$ collect the i.i.d.\ standard normal
innovations at $\vzero$. Each $Y_{k,j}$ is affine in
$(\varepsilon_{k,0},\ldots,\varepsilon_{k,j-1})$ with deterministic
coefficients, so the centred field
$G(\vartheta):=[\elln(\vartheta)-\elln(\vzero)]-\E[\cdot]$ is, for each
$\vartheta$, a centred quadratic polynomial
$G=\xi^\top M(\vartheta)\xi-\E[\xi^\top M\xi]+b(\vartheta)^\top\xi$ with
block-diagonal $M=\oplus_kM_k$. The entries of $M_k,b$ are smooth in
$\vartheta$, vanish at $\vzero$, and carry the coefficient profiles
$\gamma_k:=v^0_k/v_k-1$ (diagonal), $\delta_k\sqrt{v^0_k}\phi_{k,\cdot}/v_k$
(cross), $\delta_k^2\phi\phi'/v_k$ (drift), where
$\delta_k=\rho^0_k-\rho_k(\theta)$ and $|\phi_{k,i}|\le\sqrt{v^0_k}
(\rho^0_k)^i$, plus $u_0$-shift terms in $b$.

\emph{Step 3: second-moment and operator bounds.}
The computations of Step 2 in the proof of Theorem~\ref{thm:consistency}, keeping the
$(\sigma-\sigma_0)$-dependence, give for all $\vartheta,\vartheta'\in\Theta$:
\begin{align}
 \Var\,G(\vartheta)
 &=2\norm{M(\vartheta)}_F^2+\norm{b(\vartheta)}^2
 \ \le\ C\norm{u}^2,\label{eq:pldi-var}\\
 2\norm{M(\vartheta)-M(\vartheta')}_F^2
 +\norm{b(\vartheta)-b(\vartheta')}^2
 &\ \le\ C\norm{u-u'}^2,\label{eq:pldi-lip}\\
 \norm{M(\vartheta)}_{\mathrm{op}}\ \le\ C,
 \qquad
 \norm{M(\vartheta)-M(\vartheta')}_{\mathrm{op}}
 &\ \le\ C\norm{u-u'}\cdot\Xi_{N,n},
 \qquad \Xi_{N,n}\le C,\label{eq:pldi-op}
\end{align}
where \eqref{eq:pldi-var} collects the four component classes
($\gamma$-diagonal: $Nn(\sigma-\sigma_0)^2+T\lambda_N\Delta_n\LamN
(\theta-\theta_0)^2$; cross: $T\LamN(\theta-\theta_0)^2$; drift-quartic:
$T\LamN(\theta-\theta_0)^4\le C\,u_\theta^2\cdot\mathrm{diam}^2$;
$u_0$-terms controlled by \eqref{eq:initial-negligible}), \eqref{eq:pldi-lip}
follows from the same computations applied to the coefficient
\emph{differences}, which are products of the previous profiles with
$\lambda_k\Delta_n\LamN^{-1/2}|u_\theta-u'_\theta|$- and
$(Nn)^{-1/2}|u_\sigma-u'_\sigma|$-factors by the mean value theorem in
$\vartheta$ (each $\vartheta$-derivative of $\gamma_k,\delta_k$ produces one
factor $\lambda_k\Delta_n$ resp.\ $\sigma^{-1}$, matching the information
normalisation), and \eqref{eq:pldi-op} follows from the Schur test on the geometric row profiles:
$\norm{M_k}_{\mathrm{op}}\le C(|\gamma_k|
+\delta_k^2s_k^2/(v_k(1-\rho^0_k))
+|\delta_k|\sqrt{v^0_ks_k^2}/(v_k\sqrt{1-\rho^0_k}))\le C$ on $\Theta$,
using $1-\rho^0_k\ge c\lambda_k\Delta_n$.

\emph{Step 4: pointwise and increment tails.}
By the Hanson--Wright inequality
\citep[Theorem~1.1]{RudelsonVershynin2013} — applicable since the
$\varepsilon_{k,j}$ are independent with unit sub-Gaussian norm and
$M(\vartheta)$ is deterministic — together with the standard Gaussian tail
for $b^\top\xi$ and \eqref{eq:pldi-var}--\eqref{eq:pldi-op}:
\begin{align}
 \Pp\big(|G(\vartheta)|\ge t\big)
 &\le4\exp\Big\{-c\min\Big(\frac{t^2}{\norm u^2},\ t\Big)\Big\},
 \label{eq:pldi-hw}\\
 \Pp\big(|G(\vartheta)-G(\vartheta')|\ge t\big)
 &\le4\exp\Big\{-c\min\Big(\frac{t^2}{\norm{u-u'}^2},\
 \frac{t}{\norm{u-u'}}\Big)\Big\},\label{eq:pldi-hw-inc}
\end{align}
the increment version because $G(\vartheta)-G(\vartheta')$ is again a
Gaussian quadratic with matrix $M(\vartheta)-M(\vartheta')$ and vector
$b(\vartheta)-b(\vartheta')$, using \eqref{eq:pldi-lip}--\eqref{eq:pldi-op}.

\emph{Step 5: peeling.}
Let $\mathcal A_m:=\{\vartheta\in\Theta:2^mr\le\norm u<2^{m+1}r\}$,
$m\ge0$, $r\ge C_1$. On $\{\norm{\An(\widehat\vartheta-\vzero)}>r\}$ the
maximiser satisfies $\elln(\widehat\vartheta)\ge\elln(\vzero)$, hence by
\eqref{eq:pldi-minorant}
$\sup_{\mathcal A_m}G\ge c4^mr^2$ for some $m\ge0$. The restriction of $G$
to the ball $\{\norm u\le2^{m+1}r\}$ (radius $R_m=2^{m+1}r$) satisfies the
hypotheses of Lemma~\ref{lem:chain2d} with $u_\ast=0$ and $L=C$, by
\eqref{eq:pldi-hw-inc} (the field is a polynomial in finitely many
Gaussians, hence continuous and separable). Applying the lemma at level
$xLR_m=c4^mr^2$, i.e.\ $x=c'2^mr\ge C_1$ for $r\ge C_1/c'$:
\[
 \Pp\Big(\sup_{\mathcal A_m}G\ge c4^mr^2\Big)
 \le C\exp\{-c\,2^mr\}.
\]
Summing over $m\ge0$:
\[
 \Pp\big(\norm{\An(\widehat\vartheta_{N,n}-\vzero)}>r\big)
 \le C\,e^{-cr}\le C_Lr^{-L}
 \quad\text{for every }L,\ \text{uniformly in large }(N,n),
\]
which is \eqref{eq:PLDI}; \eqref{eq:uniform-moments} follows by tail
integration, absorbing finitely many initial $(N,n)$ into the constant via
the crude bound $\norm{\An(\widehat\vartheta-\vzero)}\le
C(\sqrt{\LamN}+\sqrt{Nn})$ on the compact $\Theta$ and the exact Gaussian
moment bounds at fixed $(N,n)$.
\end{proof}

\begin{remark}[Verification points]
Two computations deserve line-checking: the four component classes behind
\eqref{eq:pldi-var}--\eqref{eq:pldi-lip} (displayed in the proof of Theorem~\ref{thm:consistency} and differentiated here), and the constant tracking in
Lemma~\ref{lem:chain2d} (standard dyadic chaining, written self-contained
above). The Hanson--Wright input is used exactly within its hypotheses:
independent unit-variance sub-Gaussian coordinates, deterministic matrix.
\end{remark}
